% =============================================================================
% Section 10 — Comparison Analysis
% Author: David (Research Lead)
% =============================================================================

\section{Comparison Analysis}
\label{sec:comparison}

\subsection{Framing the Comparison}

Timeline Compiler is not a project-management tool, a Gantt chart engine,
or an interactive web application.
It is a \emph{visual-communication compiler}: a system that accepts structured
temporal data (or natural-language plans) and deterministically renders
presentation-quality timeline artefacts --- SVG, PNG, PDF, and PPTX --- from
a stable, version-controlled, LLM-friendly Intermediate Representation (IR).

\medskip\noindent
\textbf{What the literature says.}
The prior-art landscape reveals \emph{three} distinct clusters, not two:

\begin{itemize}
  \item \textbf{Diagram-as-code tools} (Mermaid~\cite{mermaid2023},
        PlantUML~\cite{plantuml}, D2~\cite{d2lang}, Graphviz~\cite{graphviz})
        --- source-control friendly, developer-facing, but limited in
        presentation fidelity, output breadth, and formal IR guarantees.
  \item \textbf{Declarative visualization grammars}
        (Vega-Lite~\cite{vegalite2017}, ggplot2~\cite{layeredgrammar2010},
        Vega) --- principled WHAT/HOW separation, deterministic rendering, and
        strong LLM-authorability via JSON Schema, but chart-only: they have
        no concept of node-link graphs, flows, or timeline swimlanes.
  \item \textbf{Presentation / project-management tools} (think-cell,
        PowerPoint, MS Project) --- presentation-quality, but proprietary,
        manual, and hostile to automation and version control.
\end{itemize}

\medskip\noindent
\textbf{What we recommend.}
Timeline Compiler is designed to occupy the \emph{empty cell} at the
intersection of all three clusters: open-source, git-friendly,
agent-generation-friendly, presentation-quality, \emph{and} diagram-capable
(not limited to statistical charts).
Each tool below is evaluated on the axes most relevant to that position.

% -----------------------------------------------------------------------------
\subsection{Tool-by-Tool Analysis}
% -----------------------------------------------------------------------------

\subsubsection{Mermaid (timeline and gantt)}

Mermaid~\cite{mermaid2023} is an MIT-licensed JavaScript diagramming library
with over 88{,}000 GitHub stars~\cite{mermaidwiki} and native rendering on
GitHub, GitLab, Notion, Obsidian, Azure DevOps, and hundreds of other
platforms~\cite{mermaid2023}.
Its \texttt{timeline} diagram type~\cite{mermaiddocs2024} graduated from beta
in 2023.
The syntax groups events under named sections on a horizontal time axis:

\begin{Verbatim}[fontsize=\small]
timeline
  title 2025 Platform Roadmap
  section Infrastructure
    2025 Q1 : Database migration
    2025 Q2 : CDN rollout
  section Product
    2025 Q2 : v2 public beta
    2025 Q4 : GA launch : milestone
\end{Verbatim}

\noindent
Mermaid also provides a \texttt{gantt} diagram type~\cite{mermaidgantt2024}
with sections, task durations, dependencies, and a today-marker.

\medskip\noindent
\textbf{What Mermaid does well.}
Git-friendly plain text; effortless integration into Markdown-driven
workflows; zero-install web rendering; large community; MIT license;
AI/LLM tools already generate valid Mermaid syntax.

\medskip\noindent
\textbf{Where Mermaid falls short for our goal.}
\begin{itemize}
  \item \emph{Presentation quality is limited.}
        Output is SVG/PNG rendered by a JavaScript engine. The \texttt{timeline}
        type lacks per-row swimlane formatting, precise font control, brand
        themes, or layout tuning suitable for executive slides.
  \item \emph{No PPTX or PDF output.}
        Export is browser-dependent; no native headless PPTX generation.
  \item \emph{The gantt type defaults to project-management idioms.}
        It renders task bars with dependency arrows and progress indicators ---
        appropriate for engineering sprints, not for McKinsey-style roadmaps
        where visual communication is the goal.
  \item \emph{Non-deterministic layout at scale.}
        Label placement and bar sizing change with viewport width; large
        roadmaps (>~50 events) degrade in readability.
  \item \emph{No first-class theme layer.}
        Theming is CSS injection, not a typed theme schema.
\end{itemize}

\subsubsection{PlantUML}

PlantUML~\cite{plantuml} is a GPL/commercial text-to-diagram tool built on
GraphViz.  Its \texttt{@startgantt}~\cite{plantumlgantt} mode is experimental
as of 2024 and materially less capable than Mermaid's gantt: no section
grouping, no resource assignment, no today-marker, and limited coloring.
Its timeline diagram type is likewise minimal.

PlantUML outputs PNG, SVG, ASCII art, and PDF, and integrates with
Confluence, Jenkins, and many JetBrains plugins.  The server-side rendering
architecture (a Java process) is reliable for CI pipelines.

\medskip\noindent
\textbf{Where it falls short.}
Gantt/timeline support is a secondary feature.  Syntax is verbose and
UML-centric --- not natural for non-technical stakeholders.
Visual output has a distinctly technical aesthetic; unsuitable for polished
executive deliverables.
The GPL license creates friction for proprietary embedding.

\subsubsection{vis-timeline}

vis-timeline~\cite{vistimeline} (MIT) is an interactive JavaScript timeline
widget.  Items and groups are defined as JavaScript objects; the library
renders them in the browser as scrollable, zoomable, drag-and-drop timelines.

\medskip\noindent
\textbf{Where it falls short.}
It is entirely browser-interactive: there is no declarative text format
(source data is a JavaScript data structure), no static SVG/PDF export, and
no presentation layer.  It is the right tool for dashboard UIs, not for
printed roadmaps or slide decks.  Version-controlling the rendered output
is not meaningful because the output is dynamic HTML.

\subsubsection{Frappe Gantt}

Frappe Gantt~\cite{frappegantt} (MIT) is a lightweight SVG-based Gantt
library used in ERPNext.  Tasks are provided as JavaScript objects with
\texttt{start}, \texttt{end}, \texttt{progress}, and \texttt{dependencies}
fields.  It renders to SVG in the browser.

\medskip\noindent
\textbf{Where it falls short.}
Like vis-timeline, input is programmatic JavaScript, not a text-serialisable
data format.  There is no declarative IR, no file-format input (YAML/JSON
from a file), no PPTX/PDF output, and no theme system.  The presentation
quality of the output is good for a project-tracking dashboard but lacks the
typography and layout control needed for executive presentations.

\subsubsection{Microsoft Project}

MS Project~\cite{msproject} is the industry-standard project-management
application.
It exports to a custom XML schema (\texttt{Project.xsd})~\cite{msprojectxml}
that is non-deterministic between saves: GUIDs and timestamps change on each
export, making diffs meaningless without post-processing.
Output formats are proprietary \texttt{.mpp} binary and the unstable XML.
No git-friendly text format exists natively.

\medskip\noindent
\textbf{Where it falls short.}
Proprietary, expensive (\$\textasciitilde{}30/user/month), entirely
interactive, and requires manual visual design for presentation quality.
Resource allocation and critical-path scheduling are first-class concerns ---
exactly the project-management semantics Timeline Compiler intentionally
excludes~(see \S\ref{sec:comparison-intent}).
The binary format is hostile to LLM generation.

\subsubsection{think-cell}

think-cell~\cite{thinkcell} is a PowerPoint add-in widely used by management
consultancies (McKinsey, BCG, Bain).
It produces presentation-quality Gantt/timeline charts~\cite{thinkcellgantt}
with quarter/month/week granularity, responsibility columns, and smart
label placement.
Excel data-linking enables live updates.
Pricing is approximately \$27.30/user/month~\cite{thinkcell}.

\medskip\noindent
\textbf{Where it falls short.}
Proprietary, expensive, and tightly coupled to PowerPoint.
There is no text-format input (no YAML/JSON), no CLI, no version-control
workflow, and no programmatic (agent) generation path.
Every change requires a human with PowerPoint.
Git-diffing a \texttt{.pptx} file is not meaningful.
The tool is excellent at its intended task (consulting slide production) but
solves a different problem from Timeline Compiler.

\subsubsection{PowerPoint Timeline (native SmartArt)}

PowerPoint includes SmartArt-based timeline shapes~\cite{powerpoint-timeline}.

\medskip\noindent
\textbf{Where it falls short.}
Fully manual.  No source format, no determinism, no agent path, no
version control.  Visual quality is constrained by SmartArt templates;
professional roadmaps almost always require manual box placement, which
is time-consuming and fragile when dates change.

\subsubsection{D2 (Terrastruct)}

D2~\cite{d2lang} (Mozilla Public License 2.0) is a modern Go-based
diagram language that bundles dagre, ELK~\cite{elk}, and the proprietary
TALA layout engine in a single tool.
Its key innovations are first-class container nesting, a clean
\texttt{a: label; a -> b: edge} syntax, and an official theme system
with light/dark modes~\cite{d2lang}.
D2 outputs SVG, PNG, and PDF.

\medskip\noindent
\textbf{Where it falls short for our goal.}
\begin{itemize}
  \item \emph{No timeline or Gantt diagram type.}  D2 is a general
        graph-diagram tool; swimlane roadmaps require manual
        node-positioning.
  \item \emph{TALA (the best layout engine) is closed-source and paid.}
        The open-source dagre and ELK engines produce lower-quality
        output for complex architecture diagrams.
  \item \emph{No animation.}
        D2 has a roadmap item for animation but no implementation as of 2026.
  \item \emph{No multi-panel poster layout.}
  \item \emph{MPL-2.0 copyleft implications} for embedded use.
\end{itemize}

\subsubsection{Graphviz / DOT}

Graphviz~\cite{graphviz} (Eclipse Public License) is the oldest and
most foundational text-to-diagram system, developed at AT\&T Bell Labs
circa 1991.  The DOT language declares a pure graph:
\texttt{digraph G \{ A -> B [label="step"] \}}.
Six layout engines cover different graph classes: \texttt{dot}
(Sugiyama~\cite{sugiyama1981} layered, deterministic), \texttt{neato}
(Kamada-Kawai / stress majorization~\cite{gansnerStressMaj2004},
seed-dependent), \texttt{fdp}/\texttt{sfdp}
(Fruchterman-Reingold~\cite{fruchterman1991}), \texttt{twopi}
(radial), \texttt{circo} (circular).
Output formats span SVG, PNG, PDF, PostScript, and \textasciitilde30
more~\cite{graphviz}.

\medskip\noindent
\textbf{Where it falls short for our goal.}
\begin{itemize}
  \item \emph{No timeline or presentation diagram type.}
        Graphviz is a pure graph-layout tool; swimlane roadmaps are outside
        its domain.
  \item \emph{Dated visual aesthetic.}
        The default rendering style (monospace labels, thin edges) is
        immediately recognisable as a developer tool, not an executive
        infographic.
  \item \emph{No animation and no theming system.}
  \item \emph{Force-directed engines (\texttt{neato}/\texttt{fdp}/\texttt{sfdp})
        are non-deterministic without a fixed seed.}
\end{itemize}

% -----------------------------------------------------------------------------
\subsection{Comparison Table}
\label{sec:comparison-table}
% -----------------------------------------------------------------------------

Table~\ref{tab:comparison} scores each tool on seven axes relevant to
Timeline Compiler's goals.
Ratings are qualitative (H = High / M = Medium / L = Low) based on the
analysis above and are not precise benchmarks.

\begin{table}[htbp]
\centering
\caption{Comparison of timeline and roadmap tools on axes relevant to
         Timeline Compiler.
         \textnormal{H = High, M = Medium, L = Low.
         ``Agent-gen'' = suitability for automated generation by an LLM/agent.
         ``Determinism'' = stable, reproducible output from the same input.
         ``SCM'' = source-control management friendliness.
         Viz-grammar tools scored for their supported domain (charts only).}}
\label{tab:comparison}
\small
\begin{tabular}{lccccccc}
\toprule
\textbf{Tool} &
\rotatebox{60}{\textbf{Pres.\ quality}} &
\rotatebox{60}{\textbf{Determinism}} &
\rotatebox{60}{\textbf{SCM}} &
\rotatebox{60}{\textbf{Agent-gen}} &
\rotatebox{60}{\textbf{Open source}} &
\rotatebox{60}{\textbf{PPTX out}} &
\rotatebox{60}{\textbf{SVG/PDF out}} \\
\midrule
Mermaid timeline~\cite{mermaiddocs2024}   & M & M & H & H & H & L & M \\
Mermaid gantt~\cite{mermaidgantt2024}     & L & M & H & H & H & L & M \\
PlantUML~\cite{plantuml}                  & L & H & H & M & H & L & H \\
D2~\cite{d2lang}                          & M & M & H & M & H & L & H \\
Graphviz~\cite{graphviz}                  & L & H & H & M & H & L & H \\
vis-timeline~\cite{vistimeline}           & M & L & L & L & H & L & L \\
Frappe Gantt~\cite{frappegantt}           & M & L & L & L & H & L & L \\
MS Project~\cite{msproject}               & M & L & L & L & L & L & M \\
think-cell~\cite{thinkcell}               & H & M & L & L & L & H & L \\
PowerPoint~\cite{powerpoint-timeline}     & M & L & L & L & L & H & L \\
\midrule
Vega-Lite~\cite{vegalite2017} (charts)    & H & H & H & H & H & L & H \\
ggplot2~\cite{layeredgrammar2010} (charts) & H & H & M & M & H & L & H \\
\midrule
\textbf{Timeline Compiler (target)}       & \textbf{H} & \textbf{H} & \textbf{H} & \textbf{H} & \textbf{H} & \textbf{H} & \textbf{H} \\
\bottomrule
\end{tabular}
\end{table}

\noindent
The table shows that no existing tool simultaneously achieves high scores
across all seven axes for \emph{diagram} types.
think-cell scores highest on presentation quality but is closed-source with
no agent path.
Mermaid scores highest on openness and SCM friendliness but is limited in
presentation fidelity and output format breadth.
Vega-Lite and ggplot2 prove that determinism, agent-authoring, and
presentation quality \emph{are} simultaneously achievable --- but only for
statistical charts, not for diagrams.
Timeline Compiler's design targets every axis at High --- a combination that
is, at the time of writing, unoccupied for diagram-type outputs.

% -----------------------------------------------------------------------------
\subsection{The Visualization Grammar Cluster}
\label{sec:comparison-viz-grammar}
% -----------------------------------------------------------------------------

\textbf{What the literature says.}
The grammar-of-graphics lineage --- Wilkinson's foundational
work~\cite{grammarofgraphics2005}, Wickham's operationalisation in
ggplot2~\cite{layeredgrammar2010}, and Satyanarayan et al.'s
Vega-Lite~\cite{vegalite2017} --- has conclusively demonstrated that a
small, composable, JSON-serializable grammar with a typed IR, default
inference, and a clean WHAT/HOW separation can produce presentation-quality
output, be version-controlled as plain text, and support LLM
generation~\cite{narechania2021nl4dv}.
Vega-Lite in particular is the single strongest reference architecture for
``how to design a declarative JSON grammar for
visualization''~\cite{vegalite2017}.

\medskip\noindent
\textbf{The chart/diagram gap.}
The Grammar of Graphics is explicitly scoped to statistical
graphics~\cite{grammarofgraphics2005}: it decomposes graphics into
\emph{data}, \emph{aesthetic mappings}, \emph{geometric marks},
\emph{statistical transforms}, \emph{scales}, \emph{coordinates}, and
\emph{facets}.
None of these components can express a swimlane roadmap, a node-link
architecture diagram, or a numbered step-card sequence.
Vega-Lite's \texttt{mark} vocabulary (\texttt{bar}, \texttt{line},
\texttt{point}, \texttt{rect}, etc.) does not include \texttt{node},
\texttt{edge}, \texttt{activity}, or \texttt{milestone}.
The diagram-as-code tools (Mermaid, D2, Graphviz) fill part of this gap
for developer graphs, but none delivers presentation quality,
determinism guarantees, or PPTX output.

\medskip\noindent
\textbf{What we recommend.}
The diagram compiler borrows the architectural pattern from
Vega-Lite --- (1) a small, schema-validated Domain IR, (2) a compiler that
infers defaults and validates, (3) a lower-level Scene IR, (4) multiple
rendering backends --- and extends it into the unoccupied domain of
\emph{diagram grammars}, starting with the Timeline grammar and expanding
to Flow, Graph, Comparison, and Stat-Callout grammars.
This positions the system as the ``Vega-Lite for diagrams'': the same
principled approach, applied to the visual vocabulary of technical
infographics.

% -----------------------------------------------------------------------------
\subsection{Where Timeline Compiler Intentionally Differs}
\label{sec:comparison-intent}
% -----------------------------------------------------------------------------

Three design choices explicitly diverge from the dominant approaches in
this landscape.

\paragraph{1.\ Visual-communication grammar, not task-scheduling grammar.}
Mermaid's \texttt{gantt}, MS Project, and Frappe Gantt all model a project
schedule: tasks have durations, dependencies, progress percentages, and
critical paths.
Timeline Compiler's IR models a \emph{narrative} about time: events, spans,
milestones, and swimlane groups, with an explicit visual-communication intent.
Resource allocation, dependency graphs, and critical-path analysis are
intentionally out of scope.
The Grammar of Graphics~\cite{grammarofgraphics2005} and Vega-Lite's
declarative approach~\cite{vegalite2017} inform this separation of
``what to show'' from ``how to render it''.

\paragraph{2.\ Deterministic rendering as a first-class property.}
Mermaid's layout is viewport-dependent; MS Project XML changes on every
save; PowerPoint files are binary blobs.
Timeline Compiler requires that identical IR input produces bit-identical
output.
This property is essential for CI/CD pipelines, git-based review, and
reliable LLM round-trips.

\paragraph{3.\ Agent generation as a primary use case.}
Existing tools treat human-in-the-loop as the default workflow: a user
opens an application, adjusts bars, and exports.
Timeline Compiler is designed to be driven by an AI agent with no human in
the loop: the agent generates a valid IR (JSON or YAML conforming to the
Timeline IR schema), the compiler renders it, and the artefact lands in
a pull request.
This requires a small, well-typed, documented IR that LLMs can target
reliably --- more analogous to Vega-Lite's JSON grammar~\cite{vegalite2017}
than to Mermaid's informal text syntax or MS Project's XML schema.

% -----------------------------------------------------------------------------
\subsection{Real-Data Crawl: Practical Patterns Observed}
\label{sec:comparison-patterns}
% -----------------------------------------------------------------------------

During the research phase, several real-world examples of the artefacts
Timeline Compiler targets were surveyed.
The following recurring patterns constrain the IR and rendering model.

\begin{enumerate}
  \item \textbf{Quarter granularity dominates executive roadmaps.}
        McKinsey/BCG-style roadmaps consistently use Q1--Q4 as the primary
        time unit, with sub-quarter resolution reserved for milestones.
        The IR must support a \texttt{granularity} hint (year, quarter, month,
        week) so the renderer scales the time axis appropriately.

  \item \textbf{Swimlane + milestone is the universal visual idiom.}
        Horizontal swimlanes (one per product, workstream, or business unit)
        with diamond or dot milestones at key dates is the dominant
        structural pattern across consulting firms, product teams, and
        programme offices.
        The IR needs first-class \texttt{lane} and \texttt{milestone} elements.

  \item \textbf{Spans and points coexist.}
        Real roadmaps mix spans (``Migration: Q2--Q4 2025'') and point events
        (``GA Launch: 15 Sep 2025'').
        Both must be first-class IR entities.

  \item \textbf{Azure DevOps and GitHub Projects are dominant data sources.}
        ADO work items~\cite{ado-workitems} carry iteration path
        (\texttt{System.IterationPath}) for sprint/quarter mapping.
        GitHub Projects~\cite{github-projects} use \texttt{DATE} and
        \texttt{ITERATION} custom fields for roadmap bars.
        An ingestion layer must map these fields to IR span/milestone events.

  \item \textbf{Mermaid syntax is the lowest-common-denominator input.}
        Mermaid is already widely embedded in documentation.
        An ingester that accepts Mermaid \texttt{timeline} and \texttt{gantt}
        blocks lowers the migration barrier significantly.

  \item \textbf{Manual PowerPoint is the current state-of-the-art for
        presentation quality.}
        The gap between what Mermaid can produce and what think-cell
        produces is large.
        Meeting (or approaching) think-cell output quality is the primary
        rendering challenge.
\end{enumerate}
